\chapter{Аналитическая часть}
Графовая модель -- представление кода программы в виде направленного графа.~\cite{lit1}

Основными моделями программ являются:
\begin{itemize}
	\item граф управления;
	\item информационный граф;
	\item операционная история;
	\item информационная история.
\end{itemize}
\text{}

Операционное отношение -- дуга, связывающая две вершины, когда вторая вершина может быть выполнена сразу после первой; отображает передачу управления. 

Информационное отношение -- дуга, связывающая две вершины, когда вторая вершина использует аргумент из первой; отображает передачу данных. 

Граф управления -- направленный граф, вершины которого являются операторами, а дуги -- операционными отношениями. 
 
Информационный граф -- направленный граф, вершины которого являются операторами, а дуги -- информационными отношениями. 

Операционная история -- направленный граф, вершины которого являются срабатываниями операторов, а дуги -- операционными отношениями.

Информационная история -- направленный граф, вершины которого являются срабатываниями операторов, а дуги -- операционными отношениями.

\section{Задание по варианту}

Выведите все целые числа от 1 до N включительно в порядке возрастания.

\section*{Вывод}

В данном разделе мной был проведён анализ теоретических материалов по графовым моделям.

\clearpage
